\section{All number fields from finite polynomial arithmetic}
\label{sec:dedekind}

For each monic $F\in\mathbb Z[X]$ of degree $D\ge1$ and rational prime $p$,
define an integer polynomial
\[
 U_{F,p}(T)=\prod_{g\mid\overline F}^{*}(1-T^{\deg g}).
\]
The star takes each distinct monic irreducible divisor of the reduction
$\overline F\in\mathbb F_p[X]$ once. Define $u_F(1)=1$ and, for $n>1$,
\[
 u_F(n)=\prod_{p^e\parallel n}[T^e]U_{F,p}(T),\qquad
 M_F(x)=\sum_{1\le n\le x}u_F(n).
\]
The arithmetic ERH asks for square-root growth with arbitrary positive
slack for each $M_F$, with a witness depending on $F$ and the slack.
Every monic polynomial of positive degree is admissible, including those
with repeated factors over $\mathbb Q$.

\paragraph{How one integer names the polynomial.}
Write $P(a,b)=2^a(2b+1)-1$. Each $f\ge0$ has a unique decomposition
$f=P(D-1,P(B-2,w))$ with $D\ge1$, $B\ge2$. Read the first $D$ digits
of $w$ in base $B$, starting from the units digit and filling missing
places with zeros. Decode a digit $2j$ as $j$, and a digit $2j+1$ as
$-j-1$. These are the coefficients in ascending order below the
leading $1$. Every prescribed list is represented by choosing $B$ larger
than its largest signed digit code and taking $w$ to have those digits.
In particular $f=0$ represents $X$ and $f=37$ represents $X^2+1$.

\paragraph{An exhaustive evaluation and its invariant.}
For a fixed degree $d$ there are exactly $p^d$ monic candidates $g$
over $\mathbb F_p$. Use the base-$p$ digits of $0,\ldots,p^d-1$ for their
lower coefficients. To test divisibility by a monic $g$ of degree
$d\le D$, enumerate $p^{D-d}$ monic quotients and compare the $D+1$
coefficients of the product modulo $p$. To test irreducibility, apply
that divisibility test to $g$ for all candidates of degrees $1$ through
$d-1$. A reducible monic polynomial over a field has a proper monic
divisor in this list. Thus every accepted candidate is exactly one of
the required factors, including at ramified primes.

To compute $[T^e]U_{F,p}$, start from $1$ and visit these candidates.
On accepting a degree $d$, replace the vector $v$ by
$v'_j=v_j-v_{j-d}$ for $0\le j\le e$, using zero for negative or missing
indices. After any number of visits the vector gives the product of
the accepted factors modulo $T^{e+1}$, by the coefficient rule for
multiplication. Missing high entries in the initial vector are zeros.
Trial division up to $n$ finds the contributing primes and at most $n$
divisions find each exponent. Multiplication and summation then
implement the two displayed definitions. The only special arithmetic
branch is $f=0$, where every local factor is $1-T$ and the already defined
ordinary Mertens function gives precisely the same sum.

The loops for the residue polynomials, including all quotient searches,
have the displayed primitive-recursive bounds. A vector uses at most
$e+1$ signed integers. After $j$ accepted factors its entries have
absolute value at most $2^j$, and $j\le D$, since the distinct
irreducible divisors have total degree at most $D$. These bounds also
give an explicit implementation by natural codes (for instance in base
$2^{D+2}+1$ with signed digits encoded by parity). The sum of coefficient
products for testing a divisor has at most $D+1$ terms, each at most
$p^2$. No factorization algorithm for number fields occurs in the
arithmetic formula.

\paragraph{Why the field zeta functions are recovered.}
Factor $F=\prod_{i=1}^r h_i^{m_i}$ over $\mathbb Q$. The distinct $h_i$
are monic integral polynomials by Gauss's lemma. Set
$K_i=\mathbb Q[X]/h_i$ and let $\alpha_i$ be the class of $X$.
Exclude primes dividing any nonzero
$\operatorname{disc}(h_i)$, $[\mathcal O_{K_i}:\mathbb Z[\alpha_i]]$,
or $\operatorname{Res}(h_i,h_j)$ for $i\ne j$. This excludes finitely
many primes. At all other primes the reductions of the $h_i$ are
square-free and pairwise coprime. The inclusion of orders induces
\[
 \mathbb F_p[X]/\overline h_i
       \simeq \mathcal O_{K_i}/p\mathcal O_{K_i},
\]
since the finite cokernel has order prime to $p$. By the Chinese
remainder theorem the left side splits into fields, one for each
irreducible factor $g$, with $p^{\deg g}$ elements. Hence the degrees
in $U_{F,p}$ are exactly the residue degrees of the prime ideals over
$p$ in the respective $K_i$. The powers $m_i$ disappear because the
list of factors is taken without repetitions. This is the good-prime
case of Dedekind's factorization theorem.\footnote{K.~Conrad,
\href{https://kconrad.math.uconn.edu/math5230f12/handouts/dedekindfactor.pdf}
{\emph{Factoring after Dedekind}}, Theorem 1.}

For $\sigma=\Re s>1$, absolute convergence is immediate, since
\[
 \sum_{n\ge1}|u_F(n)|n^{-\sigma}
  \le \prod_p (1+p^{-\sigma})^D \le \zeta(\sigma)^D.
\]
Indeed the absolute-coefficient sum of a local product at
$T=p^{-\sigma}$ is bounded above by the product of $1+p^{-\sigma\deg g}$.
In this half-plane comparison of the Euler factors gives
\[
 H_F(s):=\sum_{n\ge1}u_F(n)n^{-s}
     =\frac{E_F(s)}{\prod_i\zeta_{K_i}(s)}.
\]
Here $E_F$ is a finite product of factors $(1-p^{-js})^{\pm1}$
with $j\ge1$. It and its inverse are holomorphic and bounded on
$\Re s\ge b$ for each $b>0$, since
$1-p^{-jb}\le|1-p^{-js}|\le1+p^{-jb}$ there. In particular the
exceptional primes can change the values of $M_F$ but cannot mask a
zero in the critical strip.

Under ERH every $\zeta_{K_i}$ is zero-free on $\Re s>1/2$ away from
its pole at $1$ (using the standard nonvanishing on $\Re s=1$).
The reciprocal Dirichlet-series estimate proved in this report gives the
asserted growth for every coded polynomial, with its allowed
field-dependent constant. Conversely the growth bound for $M_F$ makes
$H_F$ holomorphic on that half-plane by partial summation. Given any
number field $K$, choose a primitive element $\beta$ and a positive
integer $m$ such that $m\beta$ is integral (clear denominators in a
monic rational equation). Its minimal polynomial $F$ is monic,
irreducible, represented by a natural code, and gives the single field
$K$. Now $H_F\zeta_K=E_F$ continues to the half-plane away from $1$.
The right side never vanishes there, which excludes a zero of
$\zeta_K$ to the right of the critical line. The functional equation
for the completed Dedekind zeta function, whose archimedean factors
are regular and nonzero inside the open strip, excludes a zero to the
left as well. This proves the external classical equivalence for the
entire family, including non-Galois extensions.

\paragraph{Two local computations.}
For $X^2+1$ the factors $U_{F,p}$ at $p=2,3,5$ are $1-T$, $1-T^2$,
and $(1-T)^2$ respectively. This distinguishes a repeated linear
factor, an irreducible quadratic and two split linear factors.
For $F=X^2-5$ the code at $2$ gives $1-T$ because
$\overline F=(X+1)^2$. In the maximal order of
$\mathbb Q(\sqrt5)$ the element $(1+\sqrt5)/2$ instead has minimal
polynomial $X^2-X-1$, irreducible modulo $2$, so the true inverse
Euler factor is $1-T^2$. Their ratio $1/(1+T)$ illustrates the
correction $E_F$. The finite evaluator checks the polynomial
coefficients in these examples.
